\section{Related Work}\label{sec:related}

\textbf{Chunking strategies and their selection.} Several recent studies compare
many ways of cutting documents. Shaukat et al.~\cite{shaukat2026systematic}
test thirty-six segmentation methods with five embedders and report that methods
aware of content beat plain fixed-length splitting, with paragraph groups in
first place overall. \'Smigielski et al.~\cite{smigielski2026chunking} note that
new chunkers are often built for narrow use cases and are hard to compare with
each other. Qu et al.~\cite{qu2025semantic} reach a sober view of one popular
idea: over three retrieval tasks, the gains from semantic chunking were not
steady enough to pay for its computation. Other work proposes new cutters.
HiChunk~\cite{lu2025hichunk} uses fine-tuned language models to build a
hierarchy over the document and merges units at retrieval time;
MoC~\cite{zhao2025moc} guides chunkers to output regular expressions, and
introduces two direct measures of chunk quality; StrucChunk~\cite{singh2026strucchunk}
represents legal texts as graphs of provisions and resolves cross-references
before embedding. The work closest to our main question is Adaptive
Chunking~\cite{demourajunior2026adaptive}, which picks a chunker for each document
using five measures computed from the document alone, and reports better
answers than with a single method. Its pool also contains an LLM-regex splitter and a
recursive splitter that we reuse. Orbach et al.~\cite{orbach2025hpo} tune whole
RAG configurations instead and find that greedy or random search is enough
to reach significant gains on every dataset they use.

\textbf{Retrieval units, chunk size and budgets.} How large a unit should be
has no agreed answer. Chen et al.~\cite{chen2024densex} index propositions, short units that each
hold one fact, and find better retrieval than with passages; under a given computation budget, downstream
question answering also improves. Bhat et al.~\cite{bhat2025chunksize} vary
fixed chunk sizes and find that the best one depends on the data, small for
short factual answers and large where wider context is needed, and that it also
differs between embedding models. In a chemistry setting, Amiri and
Bocklitz~\cite{amiri2025chunk} see larger effects from the embedder than from
other choices, and suggest medium to large chunks with low overlap as a starting
point. These recommendations point in different directions. Our RQ3 offers one
possible reason: in our tuning data, the preferred size changes when the number of chunks,
rather than the number of tokens, is held fixed. Whether this explains the
differences between those studies is not something we test.

\textbf{PDF parsing, OCR and benchmarks.} A second line of work asks what the
parser contributes. OmniDocBench~\cite{ouyang2024omnidocbench} evaluates document
parsing in detail over nine kinds of source.
OHR-Bench~\cite{zhang2024ocr} follows OCR noise into RAG and concludes that none
of the OCR solutions it tests is good enough for building knowledge bases. Sun
et al.~\cite{sun2026ocr} show on industrial documents that a low character error
rate does not guarantee good retrieval. Santos et al.~\cite{santos2026pdf} compare
four conversion tools on Portuguese administrative documents and attribute more
of the accuracy to metadata and hierarchy-aware splitting than to the tool.
El Bachyr et al.~\cite{elbachyr2026empirical} examine parsers and chunkers
jointly on financial questions, and Singh~\cite{singh2026parser} runs a factorial
study on Indian regulatory texts in which parser and chunker interact. Lin and
Du~\cite{lin2026pdf} retrieve layout-based evidence units from clinical
guidelines, matching page-level recall while using less context. For evaluation, UniDoc-Bench~\cite{peng2025unidoc} and
ViDoRe~V3~\cite{loison2026vidore} supply questions over visually rich pages; the
latter also has human-drawn evidence boxes, which we use as gold.

\textbf{This study.} We compare setups under a fixed token budget, and
evaluate the choice of a setup per document on two sealed held-out sets with a cross-fitted oracle,
against a decision rule written in advance. We find no evidence of a gain from
that choice in this design; studies with other data, measures or selection
signals are not directly comparable.
